% GENERATED FILE. Edit canonical lessons, then run npm run generate:books.
% canonical-source-hash: fnv1a64:e16f014dd2170a69
% canonical-lessons: TA-C261-carkkarai-noy, TA-C261-utarpayirci, TA-C261-yoka, TA-C261-arippu, TA-C261-talaiccurral

\chapter{Diabetes, Exercise, Yoga, Itch, Dizziness}
\label{ch:ta-diabetes-exercise-yoga-itch-dizziness}
\hlchaptermodality{261}

\begin{chapteropening}
\textbf{By the end of this chapter:} \emph{I can say diabetes, exercise, yoga, an itch and dizziness.}

Complete the last lesson of chapter 261: I can say diabetes, exercise, yoga, an itch and dizziness.
\end{chapteropening}

\section[\texorpdfstring{\ta{சர்க்கரை} \ta{நோய்} (carkkarai nōy)}{சர்க்கரை நோய் (carkkarai nōy)}]{\ta{சர்க்கரை} \ta{நோய்} (carkkarai nōy) — diabetes}

\label{lesson:TA-C261-carkkarai-noy}



\begin{quote}
Before the new one: say the Tamil for to float, then the Tamil for to squeeze.
\end{quote}

\subsection*{You'll want to know: \ta{சர்க்கரை} \ta{நோய்}}
\textbf{\ta{சர்க்கரை} \ta{நோய்}} — \emph{carkkarai nōy} — \textquotedblleft{}diabetes\textquotedblright{}.

\begin{grammarlens}[title={What you've built}]
One more word for this chapter.
\end{grammarlens}

\subsection*{Your turn}
\begin{itemize}
\raggedright
  \item \emph{Say it:} \emph{carkkarai nōy}
  \item \emph{Say it:} \emph{carkkarai nōy}, once more
  \item \emph{Say it:} say \emph{carkkarai nōy}
  \item \emph{Recall:} say the Tamil for to float, then the Tamil for to squeeze, then say \emph{carkkarai nōy} again
\end{itemize}

\begin{tcolorbox}[breakable,colback=teal!4,colframe=teal!35!black,title={Before you move on}]
What does \ta{சர்க்கரை} \ta{நோய்} mean? (\textquotedblleft{}Diabetes\textquotedblright{}.) Say it once more.
\end{tcolorbox}
\section[\texorpdfstring{\ta{உடற்பயிற்சி} (uṭaṟpayiṟci)}{உடற்பயிற்சி (uṭaṟpayiṟci)}]{\ta{உடற்பயிற்சி} (uṭaṟpayiṟci) — exercise}

\label{lesson:TA-C261-utarpayirci}



\begin{quote}
Before the new one: say the Tamil for to squeeze, then the Tamil for diabetes.
\end{quote}

\subsection*{You'll want to know: \ta{உடற்பயிற்சி}}
\textbf{\ta{உடற்பயிற்சி}} — \emph{uṭaṟpayiṟci} — \textquotedblleft{}exercise\textquotedblright{}.

\begin{grammarlens}[title={What you've built}]
One more word for this chapter.
\end{grammarlens}

\subsection*{Your turn}
\begin{itemize}
\raggedright
  \item \emph{Say it:} \emph{uṭaṟpayiṟci}
  \item \emph{Say it:} \emph{uṭaṟpayiṟci}, once more
  \item \emph{Say it:} say \emph{uṭaṟpayiṟci}
  \item \emph{Recall:} say the Tamil for to squeeze, then the Tamil for diabetes, then say \emph{uṭaṟpayiṟci} again
\end{itemize}

\begin{tcolorbox}[breakable,colback=teal!4,colframe=teal!35!black,title={Before you move on}]
What does \ta{உடற்பயிற்சி} mean? (\textquotedblleft{}Exercise\textquotedblright{}.) Say it once more.
\end{tcolorbox}
\section[\texorpdfstring{\ta{யோகா} (yōkā)}{யோகா (yōkā)}]{\ta{யோகா} (yōkā) — yoga}

\label{lesson:TA-C261-yoka}



\begin{quote}
Before the new one: say the Tamil for diabetes, then the Tamil for exercise.
\end{quote}

\subsection*{You'll want to know: \ta{யோகா}}
\textbf{\ta{யோகா}} — \emph{yōkā} — \textquotedblleft{}yoga\textquotedblright{}.

\begin{grammarlens}[title={What you've built}]
One more word for this chapter.
\end{grammarlens}

\subsection*{Your turn}
\begin{itemize}
\raggedright
  \item \emph{Say it:} \emph{yōkā}
  \item \emph{Say it:} \emph{yōkā}, once more
  \item \emph{Say it:} say \emph{yōkā}
  \item \emph{Recall:} say the Tamil for diabetes, then the Tamil for exercise, then say \emph{yōkā} again
\end{itemize}

\begin{tcolorbox}[breakable,colback=teal!4,colframe=teal!35!black,title={Before you move on}]
What does \ta{யோகா} mean? (\textquotedblleft{}Yoga\textquotedblright{}.) Say it once more.
\end{tcolorbox}
\section[\texorpdfstring{\ta{அரிப்பு} (arippu)}{அரிப்பு (arippu)}]{\ta{அரிப்பு} (arippu) — an itch}

\label{lesson:TA-C261-arippu}



\begin{quote}
Before the new one: say the Tamil for exercise, then the Tamil for yoga.
\end{quote}

\subsection*{You'll want to know: \ta{அரிப்பு}}
\textbf{\ta{அரிப்பு}} — \emph{arippu} — \textquotedblleft{}an itch\textquotedblright{}.

\begin{grammarlens}[title={What you've built}]
One more word for this chapter.
\end{grammarlens}

\subsection*{Your turn}
\begin{itemize}
\raggedright
  \item \emph{Say it:} \emph{arippu}
  \item \emph{Say it:} \emph{arippu}, once more
  \item \emph{Say it:} say \emph{arippu}
  \item \emph{Recall:} say the Tamil for exercise, then the Tamil for yoga, then say \emph{arippu} again
\end{itemize}

\begin{tcolorbox}[breakable,colback=teal!4,colframe=teal!35!black,title={Before you move on}]
What does \ta{அரிப்பு} mean? (\textquotedblleft{}An itch\textquotedblright{}.) Say it once more.
\end{tcolorbox}
\section[\texorpdfstring{\ta{தலைச்சுற்றல்} (talaiccuṟṟal)}{தலைச்சுற்றல் (talaiccuṟṟal)}]{\ta{தலைச்சுற்றல்} (talaiccuṟṟal) — dizziness}

\label{lesson:TA-C261-talaiccurral}



\begin{quote}
Before the new one: say the Tamil for yoga, then the Tamil for an itch.
\end{quote}

\subsection*{You'll want to know: \ta{தலைச்சுற்றல்}}
\textbf{\ta{தலைச்சுற்றல்}} — \emph{talaiccuṟṟal} — \textquotedblleft{}dizziness\textquotedblright{}.

\begin{grammarlens}[title={What you've built}]
One more word for this chapter.
\end{grammarlens}

\subsection*{Your turn}
\begin{itemize}
\raggedright
  \item \emph{Say it:} \emph{talaiccuṟṟal}
  \item \emph{Say it:} \emph{talaiccuṟṟal}, once more
  \item \emph{Say it:} say \emph{talaiccuṟṟal}
  \item \emph{Recall:} say the Tamil for yoga, then the Tamil for an itch, then say \emph{talaiccuṟṟal} again
\end{itemize}

\begin{tcolorbox}[breakable,colback=teal!4,colframe=teal!35!black,title={Before you move on}]
What does \ta{தலைச்சுற்றல்} mean? (\textquotedblleft{}Dizziness\textquotedblright{}.) Say it once more.
\end{tcolorbox}
